\chapter[逐渐而间接地获得政治意义（这种事情、我不允许）]{逐渐而间接地获得政治意义\protect\footnotemark（这种事情、我不允许）}

\footnotetext{%
    译注：``Even a purely moral act that has no hope of any immediate and visible political effect can \textbf{gradually and indirectly, over time, gain in political significance}.''
    “即使一种纯粹的道德行动并不指望产生任何直接、可见的政治效果，也可能随着时间推移，逐渐而间接地获得政治意义。”
    \par
    本章借用“道德行动获得政治意义”的句子，是为了把焦点从个人牺牲转向题材和制度层面的后果。
    小圆阻止焰、杏子试图拯救沙耶香，看似都出于私人情感，却暴露出魔法少女系统如何把善意转化为冲突。
    标题告诉读者，本章关心的是个人选择如何被更大的结构重新解释。
    所谓“不允许”，也正是在向既定规则提出伦理抗议，并开启结局前的反抗。
}

魔法少女一直都是魔女。

从外叙事层面的历史角度看，这显然没错。
魔法少女题材的出现，直接源于《家有仙妻》在一代日本女学生中意外走红\cite{ref48}。
萨曼莎是魔法少女的原型人物：
符合主流审美，呈现传统女性气质，拥有巨大的力量却被严格限制在狭小的范围之内，并受制于那种将她圈禁其中的、充满焦虑的男性气质。
她能够坚持自我，也确实常常如此，但最终还是逃不出那个时代电视行业的条条框框，被迫以完美的妆容和发型、可爱的裙装，还有一个被社会认可的妻子与母亲身份，来表演女性气质\cite{ref49}。

尽管此前在漫画中已有先例，但第一位动画魔法少女同样也是一位魔女，那就是《魔法使莎莉》中的莎莉。
她当然比《家有仙妻》里的萨曼莎年幼得多，并不承担妻子或母亲的角色，而是表演一种孩子式的女孩气质：
甜美、可爱，并且即使拥有强大的魔力，但从根本上来说，仍是无害、友善的\cite{ref50}。

这一题材便如此发展了下去。
到了《美少女战士》时，魔法少女题材的规范基本已经定型。
魔法少女像魔女一样，从异世界获得力量：
有时由镜之国的统治者赐予，有时是不小心从神秘书本中释放出来，有时则因前世曾是月之王国的公主而与生俱来。
魔法少女像魔女一样，拥有自己的使魔，也就是以动物形态出现的、有意识的生物。
当然，魔法少女也像魔女一样，可以施展强大而多样的魔法。

奇怪的是，这些力量总是指向同样来自异世界的敌人。
事实上，这些敌人往往还与她们的起源有着密切联系。
一般来说，魔法少女并不与政府腐败、企业不法行为，甚至男性英雄叙事中最常见的街头犯罪作斗争。
换言之，她们的力量不仅来自奇幻之界，而且只能用于对抗奇幻之物。

这就是所有魔法少女都是魔女的另一层含义。
魔女形象象征着对女性力量的恐惧；
在一个把男性气质等同于霸权和支配的世界里，女性气质就必须被等同于无力、服从或克制\cite{ref51}。
这一观念在日本文化中的传统表达，便是“良妻贤母”形象，也就是仁美在《魔圆》中似乎正接受训练、准备成为的那种女性理想。
良妻贤母拥有极强的社会智慧，掌握许多技艺，在家庭领域内完全主宰一切，但这些技艺都服务于家人的愉悦或支持，尤其是支持丈夫、抚养孩子。
与此同时，丈夫才是唯一被允许在家庭外行使权力的人\cite{ref52}。
日本民间文化和流行文化中充斥着“坏”女人的故事：
她们在“自己位置”之外运用力量，从原本是狐妖的妻子\cite{ref53}，到食人的山中老妖婆\cite{ref54}，再到从情人身上吸走生命温暖的雪女\footnote{%
    不过，若雪女愿意，至少最后这一点，倒也可以算作是一种慈悲了。
}\cite{ref55}。

当然，好女孩和坏女孩的二元对立绝不只局限于日本文化。
在西方民间传说和流行文化中，它的诸多表现之一，就是等待被拯救的公主，以及同时威胁公主和英雄的邪恶女巫。
掌握力量本身，就意味着成为坏女孩，成为魔女，成为对现有秩序的威胁。
\cite{ref56}

魔法少女的力量通常会通过两种方式被去威胁化。
首先，如前所述，她的力量不能触碰任何在观众眼中属于“现实”的事物，基本上只会被用来对付异世界的敌人。
其次，正如我们在第二章讨论的，她必须不断表演女性气质（记住，霸权男性气质的对立面正是表演性女性气质）：
带花边、布料少得可怜的服装、精心设计的姿势，还有变身时那近乎裸体的舞蹈，全都在提醒那些可能感受到威胁的男性观众：
她仍受制于男性凝视，遵守着“好女孩”的社会规范。

因此，沙耶香会变成魔女，直接原因就是她对两个性别歧视者使用了魔法，这绝非偶然。
她已经跨出好女孩的边界，挑战着现有秩序，因此她成了坏女孩，成了邪恶魔女。

% todo 魔圆 20:26 少革 00:08
然而，这一集通过多种方式审视并最终瓦解了那种非此即彼的对立。
最明显的一处，是在杏子与沙耶香的魔女形态 Oktavia 交手时：
蓝色与红色的血浪交织，幻化成沙耶香与杏子风格化的身影，随后又缠绕成一朵玫瑰。
这让人立刻想起《少女革命》的片头。
那部作品里也有一位“蔷薇新娘”公主，她其实是个魔女；
而她与一个持剑、拒绝传统女性角色的假小子（后者把自己定位为保护者）并肩同行，最终从顺从的“好女孩”变成强大而危险的“坏女孩”，再后来彻底挣脱了整个体系。
这一幕很可能带有普遍寓意。
因为随后那些血溅落下来，画面构图刻意让人联想到是从杏子双腿之间流出。

\begin{figure}[htbp]
    \centering
    \begin{subfigure}[t]{0.48\textwidth}
        \centering
        \includegraphics[width=\linewidth]{ch09-01.png}
        \caption{《魔法少女小圆》第九集}
        \label{fig:ch09-madoka-rose}
    \end{subfigure}
    \hfill
    \begin{subfigure}[t]{0.48\textwidth}
        \centering
        \includegraphics[width=\linewidth]{ch09-02.png}
        \caption{《少女革命》片头}
        \label{fig:ch09-utena-rose}
    \end{subfigure}
    \caption{沙耶香与杏子的血液漩涡，与《少女革命》片头中的玫瑰意象形成互文。}
    \label{fig:ch09-rose-intertext}
\end{figure}

在西方文化中，月经有时被解释为女性因天生“坏”而受到的一种特殊惩罚\cite{ref59}，因为表演性女性气质的逻辑当然就是：
只要你坚持做真实的自己，你就是“坏”，就注定要被叫做魔女。
所以每个女人骨子里都有一股“坏”劲儿，说白了，就是一个渴望表达出来的、完整而真实的自我。

可是，如果魔法少女与魔女本就是、而且从来都是一体的，那又该如何解释沙耶香的魔女化？
好在，本集给出了答案：
这是丘比强加的系统所致。
魔法少女所处的整个世界，都是由剧中唯一重要的男性角色强加给她们的。
它拥有绝对的支配权。
丘比的辩解是，她们已经“同意”参与它的系统。
但这荒唐至极，因为他蓄意隐瞒了最关键的信息。
这里明显透出一种“强奸文化”的逻辑：
处于支配地位的男性，制定出一套复杂又荒谬的标准来定义何为“同意”，然后利用这套标准把责任全部推给受害者。

但请注意，丘比实为隐含作者的能指；
这位隐含作者并不真正等同于虚渊玄，而是由编剧、导演、角色设计、动画师、配音演员、作曲家等一整个创作行业协同构成的格式塔\footnote{%
    译注：格式塔，指具有不同于部分分离特性的有机整体。
}式整体，是整个创作者产业构成的集体存在。
在这一集里，丘比明知沙耶香已无可挽救（既然她从来就是魔女，又何谈将其“从什么中拯救”呢？），却仍诱使杏子将她视为需要拯救的对象。
在成功操纵杏子接替了沙耶香曾所力图扮演的“保护者”角色之后，它又转而试图说服小圆也做出同样的牺牲。

在二次元文化中，有一个词专门指代能够唤起这种保护欲的角色（或者更准确地说，是一种角色特质）：
萌。
它源自日语“燃える”（moeru），原意为“燃烧”或“灼热”，后来一种因对方处于无助或危险境地而油然而生的、炽烈的保护欲。
2000年代后期，一种根植于此概念的审美风潮逐渐主导了整个日本二次元文化，尤其是魔法少女题材。
不过，斋藤环将其源头又再往前追溯十年，认为始于御宅族对《美少女战士》中土萌萤这一角色的反应，而这也让“萌”一词与她的名字形成了某种双关。
按照这种审美逻辑，一个角色的价值完全取决于其能否激发（假定为男性的）观众的保护欲，而能够唤起这种情感的要素包括：
无助感、“可爱”、情感脆弱以及柔弱，并且还要配上一副传统且平淡无奇的美貌\cite{ref60}。
显而易见，这与让萨曼莎屈从其懦弱丈夫的意志、迫使月野兔必须先裸身才能获得力量的逻辑如出一辙；
其本质都是让角色变得无害，从而成为“好女孩”，不致对霸权男性气质内在固有的焦虑构成任何威胁。

丘比由此被揭露为一个系统的代表：
这个系统剥夺女性的力量，强迫她们为它的利益而演出；
与此同时，又把女性置于一种处境——她们的痛苦反倒成了需要被保护的理由，从而进一步剥夺她们的力量。
换言之，它充当的是性别角色本身。
不过，它也只是一个更大、更广阔系统中的一个实例，这个系统远远超出了它本身。
《魔圆》这部作品同样也深嵌于一个更大的文化语境之中。
尽管它可以批判那些以经济可行性为名所做出的丑恶抉择，却无法彻底摆脱它们。

丘比最终被揭露为某种系统的代理人：
这一系统剥夺女性的力量，强迫她们为它的利益而演出，同时把女性置于一种处境——她们的痛苦反倒成了需要被保护的理由，从而进一步削弱了她们的力量。
换言之，丘比正是性别角色的化身。
不过，它也只是一个更庞大系统中的单一实例，该系统远不止于他自身。
《魔圆》这部作品同样也深嵌于一个更大的文化语境之中。
尽管它可以批判那些以经济可行性为名所做出的丑恶抉择，却无法彻底摆脱它们。

不过，本集的内在逻辑依然高度自洽。
沙耶香想要拯救他人、充当保护者，结果将自己异化为了魔女。
杏子和小圆试图挽救沙耶香，反而招致杀身之祸。
丘比以拯救宇宙为己任，实则延续着一个既具毁灭性又充满苦难的系统。
而正如下一集将要展现的，焰为拯救小圆所付出的一切努力，同样难逃失败的命运。

似乎，试图保护或拯救他人，从根本上就是在剥夺那个人的力量。
但另一方面，这部作品又一再批判丘比缺乏同理心，可见它并不认同安·兰德\footnote{%
    译注：Ayn Rand，本名 Alisa Zinovyevna Rosenbaum（1905年2月2日—1982年3月6日），俄裔美国作家和哲学家。
}的客观主义。
这一主义反对利他主义的强制性道德要求\cite{ref61}。
那么，也许“帮助”与“拯救”之间确有不同？
又或者，正如虚渊玄在 \emph{Fate/Zero} 后记中所写：
我们本就无能为力，一切都注定变得更糟，因为所有系统，无论宇宙、社会还是人的精神，都在无可挽回地奔向热寂\cite{ref29}？
